% Job posting: https://ca.linkedin.com/jobs/view/ai-software-developer-at-stingray-4457711563
\documentclass[11pt]{article}
\usepackage[left=1in,top=1in,right=1in,bottom=1in]{geometry}
\usepackage[hidelinks]{hyperref}
\usepackage{setspace}
\usepackage[T1]{fontenc}
\usepackage{newtxtext,newtxmath}
\ifdefined\pdfgentounicode
  \input{glyphtounicode}
  \pdfgentounicode=1
\fi
\setstretch{1.1}
\pagenumbering{gobble}
\begin{document}
\pagestyle{empty}

\noindent
Nishal Stanislaus Silva, PhD \\
Toronto, ON \\
nishal.silva@hotmail.com \, | \, +1 613 200 2030 \\
\href{https://nishal.xyz}{nishal.xyz} \, | \, \href{https://www.linkedin.com/in/nishal-silva}{linkedin.com/in/nishal-silva} \, | \, \href{https://github.com/ns2max}{github.com/ns2max}

\vspace{1.5em}
\noindent Dear Hiring Manager,

\vspace{1em}
\noindent I am writing to apply for the AI Software Developer role at Stingray. My research is built almost entirely on music-content analysis: I designed and published MIRaaS -- Music Information Retrieval as a Service -- a UDP edge-to-server offload architecture for MIR workloads, with a latency-characterization paper accepted at IEEE IS2 2026. Alongside that, I created and released four open music-content datasets on Zenodo (DoMP, DoPP, DoDP, DoDP2), totalling 9,800+ recordings from 70+ musicians. Between the two, I bring direct, concrete experience with the kind of content-analysis and catalog infrastructure a music and media company depends on.

\vspace{1em}
\noindent That research sits on top of an engineering track record in real-time audio ML: my published pattern detector runs in 14 ms on a Raspberry Pi 4 at F1 0.76, 74x faster than the DTW baseline it replaced, and I maintain Nebula, a public C++17 audio-feature-extraction and MIR toolkit with per-feature ARM latency benchmarks. I also bring 15+ years as a performing and session guitarist, which gives me musical-domain intuition that I have found genuinely changes how I approach audio-content problems, not just the engineering behind them. I already have ties to Montreal from a 2024 visiting research appointment at McGill University's IDMIL/CIRMMT lab, so the city and its music-tech community are familiar ground.

\vspace{1em}
\noindent The posting is light on day-to-day scope, so I do not yet know exactly which parts of Stingray's product and catalog stack this role touches. That is not a concern on my end -- the domain fit is exactly where I want to be -- and I would be glad to learn the specific business context and priorities quickly once I understand where the team needs the most help.

\vspace{1em}
\noindent I am a Canadian Permanent Resident and do not require sponsorship to work in Canada. I am available immediately and open to relocating to Montreal, a city I already know through my McGill collaboration. I would welcome the opportunity to discuss how my music-technology research and real-time ML engineering background could contribute to Stingray's work.

\vspace{1.5em}
\noindent Sincerely, \\
Nishal Stanislaus Silva

\end{document}
